% Appendix G -- additional tables.
\section{Additional tables}

% These numbers are recomputed from the run outputs by
% whitebox/analysis/audit.py, which fails loudly if any of them drifts.
\begin{table*}[tbp]
\caption{Where the aggregate log-probability gap comes from, synthetic control,
layer 21, $n=39$. A cell's contribution to the mean is its size share times its
own gap, so the column adds up exactly. The mean vector beats the true state
overall because F and B---the cells where the document does not work---%
contribute $133\%$ of the gap; on R, where it does work, the true state wins by
$2.5$ nats and by $0.69$ in accuracy.}
\label{tab:cells}
\centering
\small
\begin{tabular}{lrrrrr}
\toprule
cell & $n$ & mean $-$ real (nats) & contribution to $+2.747$ & acc.\ real & acc.\ mean \\
\midrule
\textbf{R} rescued    & 13 & $-2.506$ & $-0.835$ \ ($-30\%$) & \textbf{0.923} & 0.231 \\
\textbf{F} persistent & 17 & $+5.907$ & $+2.575$ \ ($+94\%$) & 0.059 & 0.294 \\
K kept                &  4 & $-0.740$ & $-0.076$ \ ($-3\%$)  & 1.000 & 0.500 \\
\textbf{B} broken     &  5 & $+8.455$ & $+1.084$ \ ($+40\%$) & 0.000 & 0.200 \\
\bottomrule
\end{tabular}
\end{table*}



\begin{table*}[tbp]
\caption{The numbers behind Figure~\ref{fig:channels}, restricted as the figure
is to the calculators whose receiver never solves an item.}
\label{tab:battery}
\centering
\small
\begin{tabular}{@{}lcc@{}}
\toprule
& accuracy & CI$_{95}$ \\
\midrule
\multicolumn{3}{@{}l}{\emph{the gold document in the prompt, and no document at all}}\\
\quad gold document in context & $0.950$ & $[0.854,\,1.000]$ \\
\quad no document & $0.050$ & $[0.000,\,0.146]$ \\
\midrule
\multicolumn{3}{@{}l}{\emph{written into the last prompt position} (best layer of 10, 14, 20, 30 for each arm)}\\
\quad replace the state with $h_{\text{skill}}$ & $0.150$ & $[0.000,\,0.306]$ \\
\quad $h+\dvec$ & $0.150$ & $[0.000,\,0.306]$ \\
\quad $h+\gvec$ & $0.050$ & $[0.000,\,0.146]$ \\
\quad $h+\dvec_j$, another instance of this skill & $0.150$ & $[0.000,\,0.306]$ \\
\quad $h+\dvec_j$, another skill & $0.100$ & $[0.000,\,0.231]$ \\
\midrule
\multicolumn{3}{@{}l}{\emph{written over the document's own span}, layer 8}\\
\quad the filler document, unpatched (the floor here) & $0.000$ & $[0.000,\,0.000]$ \\
\quad $h_{\text{recv}}+\dvec$ \ (= $h_{\text{skill}}$, $\alpha=1$) & $0.950$ & $[0.854,\,1.000]$ \\
\quad $h_{\text{recv}}+2\dvec$ & $0.800$ & $[0.625,\,0.975]$ \\
\quad $h_{\text{recv}}+\tfrac12\dvec$ & $0.050$ & $[0.000,\,0.146]$ \\
\quad $h_{\text{recv}}+\dvec$, shared prefix only & $0.900$ & $[0.769,\,1.000]$ \\
\quad $h_{\text{recv}}+\bar{\dvec}_{j}$, another skill, same prefix & $0.050$ & $[0.000,\,0.146]$ \\
\quad $h_{\text{recv}}+\dvec$, span positions permuted & $0.100$ & $[0.000,\,0.231]$ \\
\quad $h_{\text{recv}}+$ noise, $\|\cdot\|$ matched per position & $0.000$ & $[0.000,\,0.000]$ \\
\quad $h_{\text{recv}}$ itself (must be a no-op) & $0.000$ & $[0.000,\,0.000]$ \\
\bottomrule
\end{tabular}

\end{table*}

\begin{table*}[tbp]
\caption{The document-span depth sweep over all forty items, unrestricted.
Compare Table~\ref{tab:realspandepth}: the profile is the same shape with a
higher floor, because the calculators this adds are ones a wrong document
already rescues.}
\label{tab:realspan-all}
\centering
\small
\begin{tabular}{@{}rccc@{}}
\toprule
layer & accuracy & CI$_{95}$ & fraction of the document's own effect \\
\midrule
0 & $0.950$ & $[0.882,\,1.000]$ & $+1.00$ \quad\tiny(no-op arm 0.150) \\
4 & $0.950$ & $[0.882,\,1.000]$ & $+1.00$ \\
8 & $1.000$ & $[1.000,\,1.000]$ & $+1.06$ \\
12 & $0.875$ & $[0.773,\,0.977]$ & $+0.90$ \quad\tiny(no-op arm 0.150) \\
13 & $0.675$ & $[0.530,\,0.820]$ & $+0.65$ \quad\tiny(no-op arm 0.200) \\
14 & $0.700$ & $[0.558,\,0.842]$ & $+0.68$ \quad\tiny(no-op arm 0.175) \\
15 & $0.400$ & $[0.248,\,0.552]$ & $+0.29$ \quad\tiny(no-op arm 0.150) \\
16 & $0.400$ & $[0.248,\,0.552]$ & $+0.29$ \quad\tiny(no-op arm 0.175) \\
18 & $0.250$ & $[0.116,\,0.384]$ & $+0.10$ \quad\tiny(no-op arm 0.175) \\
20 & $0.475$ & $[0.320,\,0.630]$ & $+0.39$ \\
24 & $0.175$ & $[0.057,\,0.293]$ & $+0.00$ \\
30 & $0.150$ & $[0.039,\,0.261]$ & $-0.03$ \quad\tiny(no-op arm 0.150) \\
\midrule
every layer & $0.950$ & $[0.882,\,1.000]$ & $+1.00$ \\
\midrule
\emph{the filler document, unpatched} & $0.175$ & & $0.00$ \\
\emph{the gold document in context} & $0.950$ & & $1.00$ \\
\bottomrule
\end{tabular}

\end{table*}

\begin{table*}[tbp]
\caption{The same depth sweep run a second time against the other receiver---a
single fixed wrong skill rather than each skill's own paired control---on the
same items. The two agree wherever the first is well behaved; the one bump in
the first run does not replicate.}
\label{tab:realspanfixed}
\centering
\small
\begin{tabular}{@{}rccc@{}}
\toprule
layer & accuracy & CI$_{95}$ & fraction of the document's own effect \\
\midrule
0 & $0.950$ & $[0.882,\,1.000]$ & $+1.00$ \quad\tiny(no-op arm 0.125) \\
4 & $1.000$ & $[1.000,\,1.000]$ & $+1.06$ \\
8 & $0.975$ & $[0.927,\,1.000]$ & $+1.03$ \quad\tiny(no-op arm 0.125) \\
12 & $0.800$ & $[0.676,\,0.924]$ & $+0.81$ \\
16 & $0.375$ & $[0.225,\,0.525]$ & $+0.28$ \\
20 & $0.275$ & $[0.137,\,0.413]$ & $+0.16$ \quad\tiny(no-op arm 0.150) \\
\midrule
\emph{the filler document, unpatched} & $0.150$ & & $0.00$ \\
\emph{the gold document in context} & $0.950$ & & $1.00$ \\
\bottomrule
\end{tabular}

\end{table*}

\begin{table*}[tbp]
\caption{The same arms over all forty items, unrestricted, after the instrument
fix. The six calculators this adds are ones the model can partly do without the
gold document: their receiver scores $0.375$, another skill's vector $0.667$ and
norm-matched noise $0.375$, against $0.000$, $0.188$ and $0.000$ on the four
that need the document. Pooling them raises the transfer arm from $\rho=0.19$
to $0.32$, which is why the main text reports the restricted set.}
\label{tab:battery-all}
\centering
\small
\begin{tabular}{@{}lcc@{}}
\toprule
& accuracy & CI$_{95}$ \\
\midrule
\multicolumn{3}{@{}l}{\emph{the gold document in the prompt, and no document at all}}\\
\quad gold document in context & $0.950$ & $[0.854,\,1.000]$ \\
\quad no document & $0.050$ & $[0.000,\,0.146]$ \\
\midrule
\multicolumn{3}{@{}l}{\emph{written into the last prompt position} (best layer of 10, 14, 20, 30 for each arm)}\\
\quad replace the state with $h_{\text{skill}}$ & $0.150$ & $[0.000,\,0.306]$ \\
\quad $h+\dvec$ & $0.150$ & $[0.000,\,0.306]$ \\
\quad $h+\gvec$ & $0.050$ & $[0.000,\,0.146]$ \\
\quad $h+\dvec_j$, another instance of this skill & $0.150$ & $[0.000,\,0.306]$ \\
\quad $h+\dvec_j$, another skill & $0.100$ & $[0.000,\,0.231]$ \\
\midrule
\multicolumn{3}{@{}l}{\emph{written over the document's own span}, layer 8}\\
\quad the filler document, unpatched (the floor here) & $0.150$ & $[0.039,\,0.261]$ \\
\quad $h_{\text{recv}}+\dvec$ \ (= $h_{\text{skill}}$, $\alpha=1$) & $0.975$ & $[0.927,\,1.000]$ \\
\quad $h_{\text{recv}}+2\dvec$ & $0.850$ & $[0.739,\,0.961]$ \\
\quad $h_{\text{recv}}+\tfrac12\dvec$ & $0.050$ & $[0.000,\,0.118]$ \\
\quad $h_{\text{recv}}+\dvec$, shared prefix only & $0.950$ & $[0.882,\,1.000]$ \\
\quad $h_{\text{recv}}+\bar{\dvec}_{j}$, another skill, same prefix & $0.425$ & $[0.272,\,0.578]$ \\
\quad $h_{\text{recv}}+\dvec$, span positions permuted & $0.300$ & $[0.158,\,0.442]$ \\
\quad $h_{\text{recv}}+$ noise, $\|\cdot\|$ matched per position & $0.200$ & $[0.076,\,0.324]$ \\
\quad $h_{\text{recv}}$ itself (must be a no-op) & $0.125$ & $[0.023,\,0.227]$ \\
\bottomrule
\end{tabular}

\end{table*}

\begin{table*}[tbp]
\caption{What the content vector contains once it is able to contain the
task---the numbers behind Figure~\ref{fig:skilltask}.}
\label{tab:skilltask}
\centering
\small
\begin{tabular}{@{}lccc@{}}
\toprule
donor written over the document's span & accuracy & CI$_{95}$ & fraction \\
\midrule
\emph{the filler document, unpatched} & $0.000$ & & $0.00$ \\
\emph{the gold document in context} & $0.850$ & & $1.00$ \\
\midrule
$\dvec_i$, this instance's own & $0.850$ & $[0.694,\,1.000]$ & $+1.00$ \\
$\bar{\dvec}_{\neg i}$, the mean over the skill's other instances & $0.800$ & $[0.625,\,0.975]$ & $+0.94$ \\
$\dvec_j$, one other instance & $0.900$ & $[0.769,\,1.000]$ & $+1.06$ \\
$\dvec_i^{\parallel}$, the half along that mean & $0.850$ & $[0.694,\,1.000]$ & $+1.00$ \\
$\dvec_i^{\perp}$, the rest & $0.000$ & $[0.000,\,0.000]$ & $+0.00$ \\
$\bar{\dvec}$, another skill & $0.150$ & $[0.000,\,0.306]$ & $+0.18$ \\
$\bar{\dvec}$, averaged over every skill & $0.050$ & $[0.000,\,0.146]$ & $+0.06$ \\
the same $\dvec_i$, span positions permuted & $0.000$ & $[0.000,\,0.000]$ & $+0.00$ \\
noise, norm matched per position & $0.000$ & $[0.000,\,0.000]$ & $+0.00$ \\
\bottomrule
\end{tabular}

\end{table*}


\label{app:tables}

\begin{table*}[tbp]
\caption{Three transplant channels, synthetic control, Qwen3-1.7B, $n=39$.
Each row is a separate experiment; the receiver baseline differs because the
document-span and question-span transplants run inside a prompt that still
contains a filler document.}
\label{tab:windows}
\centering
\small
\begin{tabular}{llccc}
\toprule
channel & positions transplanted & effective layers & peak acc. & receiver baseline \\
\midrule
document span & the document's own 688 tokens & 0--10 & 0.436 (L3--4) & 0.154 \\
question span & the question's 42--47 tokens  & 11--16 & 0.333 (L12/14) & 0.154 \\
final position & the last prompt token        & 21--27 & 0.436--0.487 & 0.231 \\
\bottomrule
\end{tabular}
\end{table*}


\begin{table*}[tbp]
\caption{The same measurements at two model sizes, on the $358$-item reruns.
The effect itself, in both formats, is the top block of Table~\ref{tab:format}.
The multiple-choice presence row is the reason we take the free-form format as
the primary evidence: on the 8B, chance on four options is $0.250$ and $\gvec$
sits at $0.298$, within a few points of it, while in free-form it reads
$0.008$.}
\label{tab:ladder}
\centering
\small
\begin{tabular}{lcc}
\toprule
 & Qwen3-1.7B & Qwen3-8B \\
\midrule
MC, rescued cell: replace outright        & 1.000 & 1.000 \\
MC, rescued cell: $+\dvec$                & 0.805 & 0.894 \\
MC, rescued cell: $+\gvec$                & 0.098 & 0.298 \\
free-form, rescued cell: replace outright & ---   & 0.235 \\
free-form, rescued cell: $+\dvec$         & ---   & 0.210 \\
free-form, rescued cell: $+\gvec$         & ---   & 0.008 \\
\midrule
layers where $\|\gvec\| > \|\tvec\|$      & 0--16 & 0--12 \\
$\mathrm{PR}(\dvec)$ centred, deepest layer & 5.4 & 18.1 \\
$\cos(\dvec_i,\dvec_j)$ same / other family, deepest & 0.869 / 0.705 & 0.584 / 0.530 \\
\bottomrule
\end{tabular}
\end{table*}

\begin{table*}[tbp]
\caption{The content component split along its own shared direction and
orthogonal to it, rescued cell, $n_R=82$. The crossover at layer 21 is the same
layer at which cumulative attention knockout on the document becomes free
(\S\ref{sec:where}) and at which another item's content vector stops working
(\S\ref{sec:whose}).}
\label{tab:parperp}
\centering
\small
\begin{tabular}{lcccccc}
\toprule
injected & L16 & L18 & L20 & L21 & L24 & L26 \\
\midrule
$h_{\text{no}} + \dvec_i$ (whole) & 0.134 & 0.390 & 0.439 & 0.500 & 0.695 & 0.805 \\
$h_{\text{no}} + (\dvec_i\!\cdot\!u)u$ (shared dir.) & 0.122 & 0.305 & 0.268 & 0.098 & 0.049 & 0.085 \\
$h_{\text{no}} + \dvec_i^{\perp}$ (item-specific) & 0.012 & 0.122 & 0.195 & 0.427 & 0.524 & 0.561 \\
$h_{\text{no}} + \bar{\dvec}$ (shared dir., fixed norm) & 0.134 & 0.305 & 0.293 & 0.122 & 0.049 & 0.085 \\
\bottomrule
\end{tabular}

\end{table*}

\begin{table*}[tbp]
\caption{The decomposition of Eq.~\ref{eq:decomp}, synthetic control,
Qwen3-1.7B, fp32, $n=358$, at the last prompt position. $\mathrm{PR}$ is the
participation ratio of $\{\dvec_i\}$ after removing their shared mean direction:
the effective number of dimensions the item-to-item variation occupies, out of
2048. The last column is the fraction of a single item's $\dvec_i$ that lies
along that shared direction.}
\label{tab:geometry-tierA}
\centering
\small
\begin{tabular}{rrrrrrrrr}
\toprule
layer & $\|\dvec\|$ & $\|\gvec\|$ & $\|\dvec\|/\|\tvec\|$ & $\cos(\dvec,\gvec)$
& \multicolumn{2}{c}{$\cos(\dvec_i,\dvec_j)$} & $\mathrm{PR}(\dvec)$ & on mean \\
\cmidrule(lr){6-7}
& & & & & same fam. & cross fam. & & \\
\midrule
 0 &     1.01 &     2.66 & 0.43 & $-0.50$ & 0.996 & 0.995 & 3.2 & 0.998 \\
10 &    16.24 &    38.96 & 0.45 & $-0.38$ & 0.989 & 0.987 & 4.5 & 0.994 \\
14 &    33.47 &    68.53 & 0.52 & $-0.37$ & 0.979 & 0.970 & 8.6 & 0.987 \\
21 &   459.14 &   352.60 & 1.04 & $-0.43$ & 0.920 & 0.794 & 5.5 & 0.919 \\
27 &   919.82 &   704.45 & 1.00 & $-0.37$ & 0.869 & 0.705 & 5.4 & 0.878 \\
\bottomrule
\end{tabular}

\end{table*}
\begin{table*}[tbp]
\caption{Fraction of the rescued cell answered correctly under each injected
component. Synthetic control, multiple choice, Qwen3-1.7B, fp32, $n=358$,
rescued cell $n_R=82$. Bootstrap 95\% CIs over items are $\pm0.09$ to $\pm0.11$
at these magnitudes and are given in the appendix. $\bar{\dvec}$ is the mean
content direction over items, restored to the typical per-item norm; $\dvec_j$
is another item's content vector, drawn from the same or a different conversion
family. The row that licenses the name ``content component'' is $\gvec$.}
\label{tab:causal}
\centering
\small
\begin{tabular}{lccccc}
\toprule
injected at the last prompt position & L17--20 & L21 & L24 & L26 & L27\\
\midrule
$h_{\text{skill}}$ (replace outright; upper ref.) & 0.29--0.45 & 0.988 & 0.963 & 1.000 & 1.000 \\
$h_{\text{skill}}$ into the wrong-document prompt & 0.51--0.77 & 0.976 & 0.963 & 1.000 & 1.000 \\
\midrule
$h_{\text{no}} + 0.5\,\dvec_i$ & 0.15--0.27 & 0.232 & 0.280 & 0.317 & 0.268 \\
$h_{\text{no}} + \dvec_i$ \quad (content) & 0.32--0.44 & 0.500 & 0.695 & 0.805 & 0.780 \\
$h_{\text{no}} + 2\,\dvec_i$ & 0.43--0.67 & 0.841 & 0.976 & 0.976 & 0.951 \\
\midrule
$h_{\text{no}} + \gvec_i$ \quad (presence) & 0.00--0.04 & 0.098 & 0.098 & 0.098 & 0.098 \\
$h_{\text{no}} + \bar{\dvec}$ \ (one shared direction) & 0.24--0.32 & 0.122 & 0.049 & 0.085 & 0.098 \\
$h_{\text{no}} + \dvec_j$ \ (same family) & 0.29--0.37 & 0.268 & 0.244 & 0.317 & 0.293 \\
$h_{\text{no}} + \dvec_j$ \ (other family) & 0.18--0.23 & 0.073 & 0.049 & 0.049 & 0.049 \\
\bottomrule
\end{tabular}

\end{table*}
